\documentclass[../DoAn.tex]{subfiles}
\begin{document}

\section{Kết luận}

Trong phạm vi đề tài, đồ án đã xây dựng được một hệ thống mô phỏng và hiển thị
giao thông trong không gian ba chiều, trong đó \gls{SUMO} đảm nhiệm mô phỏng vi
mô, Python điều khiển mô phỏng theo từng bước thời gian thông qua \gls{TraCI}, và
Unity đảm nhiệm dựng cảnh, hiển thị theo thời gian thực và xử lý tương tác. So
với cách quan sát kết quả mô phỏng truyền thống vốn chủ yếu dựa trên góc nhìn 2D
từ trên cao, hệ thống của đồ án không chỉ tự động hóa luồng dữ liệu và đồng bộ
theo bước mô phỏng mà còn cho phép người dùng trải nghiệm theo góc nhìn của người
tham gia giao thông, qua đó hỗ trợ quan sát và phân tích trực quan hơn trong các
kịch bản có mật độ đối tượng cao hoặc khi cần thuyết minh kết quả.

Về kết quả đạt được, đồ án đã hoàn thành các nhóm chức năng chính sau:

\begin{itemize}[leftmargin=2em]
	\item \textbf{Kết nối và đồng bộ thời gian thực}: triển khai kênh giao tiếp
	giữa Unity và Python qua TCP cùng cơ chế tách khung bản tin bằng dấu hiệu kết
	thúc \texttt{<END>} để khắc phục đặc thù truyền theo luồng của TCP; hỗ trợ
	đồng thời bản tin dữ liệu trạng thái và bản tin điều khiển.
	\item \textbf{Tự động dựng cảnh từ dữ liệu mạng đường}: dựng mạng đường trong
	không gian ba chiều từ kịch bản sinh tự động theo lưới mê cung phục vụ kiểm
	thử, hoặc từ bản đồ thực tế OpenStreetMap, giảm đáng kể thao tác thủ công khi
	thay đổi kịch bản.
	\item \textbf{Hiển thị đa đối tượng theo thời gian}: tạo, cập nhật và thu hồi
	phương tiện, đèn tín hiệu và người đi bộ theo vòng đời mô phỏng; áp dụng nội
	suy chuyển động để hiển thị mượt khi tốc độ khung hình cao.
	\item \textbf{Tương tác lái xe}: cho phép người dùng chiếm quyền điều khiển
	một phương tiện và lái thủ công bằng cơ chế vật lý; xử lý va chạm để chuyển
	xe bị tông thành trạng thái ``xác xe'' gây ùn ứ; và trả quyền điều khiển kèm
	cơ chế đưa xe trở lại làn đường gần nhất để mô phỏng tiếp tục, hoặc loại bỏ xe
	khi không thể đưa về làn.
	\item \textbf{Hai chế độ vận hành và lưu trữ phiên}: hỗ trợ chế độ thời gian
	thực và chế độ phát lại; gom toàn bộ kết quả của mỗi lần chạy vào một thư mục
	phiên thống nhất (nhật ký hành trình, dữ liệu mạng, chuỗi trạng thái theo bước
	và bản tóm tắt), phục vụ truy vết và phát lại.
	\item \textbf{Điều khiển vận hành}: cung cấp các thao tác điều khiển camera,
	tạm dừng và tiếp tục, điều chỉnh tốc độ mô phỏng lúc đang chạy, và lọc đối
	tượng hiển thị theo vùng quan sát nhằm duy trì hiệu năng.
\end{itemize}

Bên cạnh các kết quả đạt được, hệ thống vẫn còn một số hạn chế cần tiếp tục hoàn
thiện. Thứ nhất, giao thức truyền thông hiện tại ưu tiên tính đơn giản nên khi số
lượng đối tượng lớn có thể phát sinh tải xử lý và băng thông, làm giảm tốc độ cập
nhật. Thứ hai, mô hình vật lý cho chế độ lái còn ở mức cơ bản, chưa phản ánh đầy
đủ động lực học của phương tiện; cơ chế đưa xe trở lại làn đường sử dụng ngưỡng
khoảng cách cố định nên trong một số tình huống phức tạp có thể chưa tối ưu. Thứ
ba, hệ thống hiện chạy trên một máy và hỗ trợ một người dùng cho mỗi phiên mô
phỏng, chưa hỗ trợ nhiều người dùng cùng tham gia một phiên.

Từ quá trình thực hiện, đồ án rút ra một số kinh nghiệm chính: (i) trong hệ thống
thời gian thực, cần thiết kế cơ chế phân định ranh giới thông điệp rõ ràng khi
dùng TCP để tránh các lỗi khó truy vết; (ii) việc chuẩn hóa dữ liệu và thống nhất
hệ tọa độ, đơn vị ngay từ đầu giúp giảm đáng kể công sức tích hợp giữa mô phỏng
và hiển thị; và (iii) khi xe được điều khiển theo nhiều trạng thái khác nhau,
việc quản lý vòng đời và quyền điều khiển một cách tập trung là then chốt để hệ
thống vận hành ổn định.

\section{Hướng phát triển}

Trong thời gian tới, đồ án có thể được phát triển theo các hướng chính sau:

\begin{itemize}[leftmargin=2em]
	\item \textbf{Tối ưu truyền thông và hiệu năng}: giảm chi phí đóng gói và phân
	tích dữ liệu bằng định dạng nhị phân hoặc nén, hoặc chỉ truyền phần thay đổi
	theo từng bước; bổ sung cơ chế đo độ trễ đầu--cuối để đánh giá chất lượng đồng
	bộ khi số lượng đối tượng tăng.
	\item \textbf{Hỗ trợ nhiều người dùng đồng thời}: mở rộng kiến trúc để nhiều
	máy hiển thị cùng tham gia một phiên mô phỏng, mỗi người điều khiển một phương
	tiện riêng; đây là tính năng chưa có trong phiên bản hiện tại và đòi hỏi cơ
	chế đồng bộ quyền điều khiển và trạng thái giữa các máy.
	\item \textbf{Nâng cấp mô hình lái và xử lý va chạm}: cải thiện động lực học
	phương tiện cho chế độ lái thủ công; tinh chỉnh cơ chế đưa xe trở lại làn
	đường theo ngữ cảnh thay vì ngưỡng khoảng cách cố định, giúp xe hòa lại dòng
	giao thông tự nhiên hơn.
	\item \textbf{Mở rộng kịch bản và đánh giá định lượng}: xây dựng bộ kịch bản
	chuẩn theo nhiều kích thước mạng và mật độ phương tiện; chuẩn hóa các thước đo
	như tốc độ khung hình, độ trễ và thời gian di chuyển, đồng thời tự động hóa
	quá trình chạy thực nghiệm để có cơ sở so sánh khách quan.
	\item \textbf{Cải thiện chất lượng phần mềm}: bổ sung kiểm thử cho các mô-đun
	sinh mạng, sinh tuyến và đóng gói giao thức; chuẩn hóa cấu hình chạy và tài
	liệu hóa giao thức nhằm thuận tiện cho việc bảo trì và mở rộng về sau.
\end{itemize}

\end{document}
